\begin{figure*}[t]
\centering
\includegraphics[width=0.84\textwidth]{validation-paired-outcomes.pdf}
\caption{Paired outcomes and recorded trigger burden. (a) Full-minus-comparator
differences with 95\% participant-bootstrap intervals; sample sizes appear beside
each endpoint. (b) Participant rates over 29.164 reference-negative online hours.
Lines connect the same participant; fixed horizontal offsets separate overlapping
participants without altering rates.
Rates describe the enumerated event set.}
\label{fig:participant_outcomes}
\end{figure*}

\subsection{Scripted Checks, Simulated Signs, and Natural Use}
\label{sec:participant_validation}
We recorded 42 sessions from 30 participants between 17 August and 5 September
2026. Session categories described the acquisition setting, not the event label.
Six \textit{scripted-action} sessions involved routine instructed B/E/F/A/S checks
without target symptoms. Seven \textit{simulated-positive} sessions included
volunteers enacting facial droop, arm drift, trunk lean, or altered articulation.
Five \textit{challenge} sessions combined simulated signs with difficult inputs,
including occlusion, low light, background sound, incomplete actions, and disconnection.
Twelve natural-use sessions measured ordinary-activity triggering; 12 repeat
sessions assessed repeated measurements. The 38 annotated targets comprised
28 events from simulated-positive sessions and 10 from challenge sessions.

The recorded configuration, v1.3.0+frozen (Full), used a personal Balance baseline
and quality gates. Two historical comparison labels denoted a population Balance reference
(Population baseline) or disabled gates (No gates). All three shared 201 event
IDs. Annotations were recorded as blinded to system outputs. We analysed paired
event outcomes with 10,000 participant-cluster bootstrap resamples (seed 20260910)
and pointwise 95\% intervals. The seed fixes resampling for reproducibility.
Target response counted the recorded trigger label, without inferring passive
detection or completed guided assessment. All 38 runtime paths remain unresolved;
component names and acquisition timestamps do not establish their origin.
Completeness indicated a completed
measurement process linked to a report, independently of positivity.

Full had recorded triggers for 35/38 targets (92.1\%) from 12 participants, versus 29/38 for
Population baseline and 33/38 for No gates (Fig.~\ref{fig:participant_outcomes}).
Face, Arms, Balance, and Speech counts were 9/10, 7/8, 10/11, and 9/9.
The 15.8-percentage-point difference from Population baseline (95\% interval
7.9--23.7) came entirely from six additional Balance detections.
The difference from No gates was 5.3 points (0.0--14.7).
Exact population-reference parameters and disabled rules remain unspecified,
so these are descriptive contrasts between historical labels, not controlled
ablations of the current implementation.

Among 177 started measurements from 25 participants, 159 were reference-complete
and 18 were incomplete, including two aborted attempts. All configurations
retained the 159 complete inputs. Incorrect completions were 1/18 for Full,
1/18 for Population baseline, and 15/18 for No gates. The Full-minus-No gates
difference was $-77.8$ points ($-95.0$ to $-55.6$), describing a recorded completion difference. The unavailable gate configuration
precludes attributing it to a specific quality rule.

During 29.164 reference-negative online hours from 12 natural-use participants,
Full, Population baseline, and No gates recorded 14, 20, and 19 false triggers
(0.480, 0.686, and 0.651 per hour). Online availability was 92.8\%.
These rates describe recorded triggers; complete capture coverage and linked
check outcomes are needed to quantify final-alert burden. Full also retained
two symptom-count shortfalls among 35 known target records and one incomplete
correction history among nine cases. Three target symptom counts were unknown.
All three Full urgency shortfalls occurred among missed targets; detected targets
had no urgency-label mismatches. Supporting notes retain extended endpoint and timing definitions.

Face reference values were recorded as the mean of two independent geometric
annotations, with uncertainty listed as $\pm0.006$ but no statistical definition.
The exported system values differed from these references by a mean absolute error (MAE)
of 0.00452 interocular-width units (33 measurements, 22 participants;
Fig.~\ref{fig:participant_reference}). The difference is comparable to the
listed reference uncertainty and does not establish measurement accuracy.
Eleven complete repeat pairs, separated
by 11.2--15.6 days, yielded ICC(2,1)=0.935 (95\% interval 0.801--0.975);
one participant lacked the first measurement. Eyes exported gaze offset had
MAE 0.280$^\circ$ against desktop eye-tracker values (25 measurements,
19 participants; system range 0.094--1.641$^\circ$). Its runtime output mapping
remains unresolved, so this comparison does not establish angular accuracy.
Automatic Eyes abnormality classification remained disabled.

\begin{figure*}[t]
\centering
\includegraphics[width=0.84\textwidth]{validation-reference-repeatability.pdf}
\caption{Exported-value comparisons and Face repeatability.
(a,b) System-minus-reference errors for Face mouth-corner change and Eyes gaze
offset; dashed lines indicate zero difference. Eyes compares exported fields with
unresolved runtime mapping. (c) Two system measurements for each
of 11 participants with complete Face repeat pairs. Lines connect participants
across visits.}
\label{fig:participant_reference}
\end{figure*}
